\section{Dari Ekspresi Reguler ke Automata dan Token Recognition}
\label{sec:ch06-thompson}

\subsection{Konstruksi Thompson}

Konstruksi Thompson mengubah ekspresi reguler menjadi NFA dengan
transisi epsilon. Setiap ekspresi reguler memiliki NFA dengan satu
state awal dan satu state penerima. Aturan dasarnya:
\begin{table}[H]
\centering
\caption{Aturan konstruksi Thompson}
\label{tab:ch06-thompson}
\begin{tabularx}{\textwidth}{@{}l X@{}}
\toprule
\textbf{Ekspresi} & \textbf{Konstruksi} \\
\midrule
\(\varnothing\) & State awal dan penerima tanpa jalur \\
\(\varepsilon\) & Transisi epsilon dari awal ke penerima \\
\(a\) & Transisi \(a\) dari awal ke penerima \\
\(r \mid s\) & State awal baru bercabang ke NFA \(r\) dan \(s\) \\
\(rs\) & State penerima \(r\) dihubungkan ke awal \(s\) \\
\(r^{*}\) & Loop epsilon di sekitar NFA \(r\) \\
\bottomrule
\end{tabularx}
\end{table}

NFA hasil konstruksi mungkin memiliki banyak transisi epsilon. Setelah
itu, subset construction menghasilkan DFA, dan minimalisasi
menghilangkan state yang ekuivalen. Alur lengkapnya adalah:
\[
\text{regex} \to \text{NFA Thompson} \to \text{DFA} \to
\text{DFA minimal}.
\]

\subsection{Epsilon Closure dalam Kode}

Epsilon closure adalah inti simulasi NFA dan subset construction.
Implementasinya menggunakan penelusuran graf. Potongan kode Python
ditunjukkan pada \lstref{lst:ch06-epsilon}.

\includecode[
  language=Python,
  firstline=127,
  lastline=136,
  firstnumber=127,
  caption={Fungsi \code{epsilon\_closure} pada \file{fa\_tool.py}},
  label={lst:ch06-epsilon}
]{../../code/tools/fa_tool.py}

Fungsi ini mengembalikan \code{frozenset} agar dapat dipakai sebagai
kunci kamus transisi DFA. Penggunaan himpunan beku juga mencegah
modifikasi tidak sengaja.

\subsection{Menggabungkan Beberapa Token}

Lexer biasanya mengenali banyak token sekaligus. Caranya adalah
membuat NFA terpisah untuk setiap pattern, lalu menambahkan state awal
baru dengan transisi epsilon ke setiap NFA. State penerima setiap NFA
diberi label jenis token. Jika beberapa token cocok pada panjang yang
sama, label dengan prioritas tertinggi yang dipilih.

\begin{figure}[H]
\centering
\begin{tikzpicture}[node distance=6mm and 10mm]
  \node[state, initial] (start) {\(s\)};
  \node[phase, text width=2.6cm, right=of start] (id) {NFA\\identifier};
  \node[phase, text width=2.6cm, below=of id] (num) {NFA\\angka};
  \node[phase, text width=2.6cm, below=of num] (op) {NFA\\operator};
  \draw[arrow] (start) -- node[above, font=\small] {\(\varepsilon\)} (id);
  \draw[arrow] (start) -- node[above, font=\small] {\(\varepsilon\)} (num);
  \draw[arrow] (start) -- node[above, font=\small] {\(\varepsilon\)} (op);
\end{tikzpicture}
\caption{Menggabungkan NFA token dengan state awal baru}
\label{fig:ch06-gabung}
\end{figure}

\subsection{Maximal Munch pada DFA}

DFA gabungan mengenali seluruh token. Untuk menerapkan maximal munch,
lexer menjalankan DFA selama masih ada transisi yang mungkin, sambil
mengingat posisi terakhir yang berada di state penerima. Ketika
transisi berikutnya tidak ada, lexer kembali ke posisi terakhir
tersebut dan mengeluarkan token. Sisa karakter diproses pada
pemanggilan berikutnya.

Pendekatan ini memerlukan \eng{rollback} pada buffer masukan. Karena
itu, lexer biasanya menyimpan seluruh berkas atau blok besar di
memori, bukan membaca karakter satu per satu dari disk.

\subsection{Dari Otomata ke Lexer}

DFA gabungan dapat direpresentasikan sebagai tabel transisi
dua dimensi: baris adalah state, kolom adalah simbol, dan isi sel
adalah state tujuan. Tabel ini dapat dikompilasi menjadi array atau
\eng{switch} besar. Lexer tangan biasanya memakai \eng{switch}
untuk keterbacaan, sedangkan lexer generator memakai tabel untuk
kecepatan.

Bab VII akan mengimplementasikan kedua pendekatan: lexer tangan
dalam C dan lexer Flex. Perbandingan keduanya menunjukkan trade-off
antara kontrol dan produktivitas.
